\begin{tikzpicture}[scale=0.92,>=Stealth]
  \tikzset{
    rope/.style={thick},
    inc/.style={blue,very thick,->},
    ref/.style={orange!90!black,very thick,->},
    wall/.style={black,very thick},
    note/.style={font=\small,align=center}
  }

  \begin{scope}
    \node[font=\bfseries] at (0,1.95) {固定端};
    \draw[rope] (-3.2,0) -- (2.6,0);
    \draw[wall] (2.6,-1.05) -- (2.6,1.05);
    \foreach \y in {-0.9,-0.55,...,0.9}{
      \draw (2.6,\y) -- (2.92,\y+0.18);
    }
    \draw[blue,thick,samples=80,domain=-2.90:-0.55]
      plot(\x,{0.72*exp(-2.4*(\x+1.75)^2)});
    \draw[inc] (-2.55,0.88) -- (-1.35,0.88) node[midway,above,font=\scriptsize] {入射脉冲};
    \draw[orange!90!black,thick,samples=80,domain=0.05:2.25]
      plot(\x,{-0.62*exp(-2.4*(\x-1.15)^2)});
    \draw[ref] (2.10,-0.86) -- (0.95,-0.86) node[midway,below,font=\scriptsize] {反射脉冲};
    \node[note] at (0,-1.65) {反射后倒相\\端点位移始终为 $0$，形成波节};
  \end{scope}

  \begin{scope}[xshift=7.10cm]
    \node[font=\bfseries] at (0,1.95) {自由端};
    \draw[rope] (-3.2,0) -- (2.6,0);
    \draw[thick] (2.6,-0.10) circle (0.16);
    \draw[thick] (2.76,-0.10) -- (3.05,-0.10);
    \draw[blue,thick,samples=80,domain=-2.90:-0.55]
      plot(\x,{0.72*exp(-2.4*(\x+1.75)^2)});
    \draw[inc] (-2.55,0.88) -- (-1.35,0.88) node[midway,above,font=\scriptsize] {入射脉冲};
    \draw[orange!90!black,thick,samples=80,domain=0.05:2.25]
      plot(\x,{0.62*exp(-2.4*(\x-1.15)^2)});
    \draw[ref] (2.10,0.86) -- (0.95,0.86) node[midway,above,font=\scriptsize] {反射脉冲};
    \node[note] at (0,-1.65) {反射后不倒相\\端点振幅最大，形成波腹};
  \end{scope}

  \node[note] at (3.55,-2.38)
    {边界条件决定反射相位：固定端是波节，自由端是波腹};
\end{tikzpicture}
